\documentclass[11pt]{article}
\usepackage{mathblatt}
% gesetzt von werkzeuge/setzer.py (Sorte selbst) aus dem Lernweg VUD
% und den Bankzeilen; Muster bau/proben/2026-10-09/hypotenuse-muster.
\usepackage{enumitem}
\usepackage{adjustbox,varwidth}
\newgeometry{a4paper,margin=18mm,top=15mm,bottom=20mm,footskip=9mm}
\pagestyle{fancy}\fancyhf{}\renewcommand{\headrulewidth}{0pt}
\fancyfoot[L]{\footnotesize\color{mbgrau}VUD \textperiodcentered{} Potenzen}
\fancyfoot[R]{\footnotesize\color{mbgrau}\thepage}
\setlength{\parskip}{3pt}
\renewcommand{\kreuz}[1]{\mbox{$\square$\ #1}\quad}
\newcommand{\kreuzl}[1]{\par\hangindent1.4em\hangafter1\noindent$\square$\ #1}
\newcommand{\aufg}[3]{\par\noindent\makebox[0pt][r]{#3}\begin{minipage}[t]{\linewidth}#1\end{minipage}\par}
\newcommand{\aufgg}[4]{\par\noindent\makebox[0pt][r]{#4}\begin{minipage}[t]{0.6\linewidth}\vspace{0pt}#1\end{minipage}\hfill
  \begin{minipage}[t]{0.37\linewidth}\vspace{0pt}\centering #2\end{minipage}\par}
\newcommand{\nr}[1]{\makebox[7mm][l]{\textbf{#1}}}
\newcommand{\lsg}[3]{\par\noindent\begin{minipage}[t]{9mm}\textbf{#1}\end{minipage}%
  \begin{minipage}[t]{52mm}\raggedright\bfseries\boldmath #2\end{minipage}\hspace{3mm}%
  \begin{minipage}[t]{\dimexpr\linewidth-64mm\relax}\raggedright\small #3\end{minipage}\par\vspace{4pt}}
\newlength{\kr}
\makeatletter
\newcommand{\karomerk}[2]{\protected@write\@auxout{}{\string\karoseite{#1}{#2}{\thepage}}}
\newcommand{\karoseite}[3]{\expandafter\xdef\csname karo@#1@#2\endcsname{#3}}
\newcommand{\karohinweis}[3]{\karomerk{h}{#1}\ifcsname karo@k@#1\endcsname
  \ifnum\csname karo@k@#1\endcsname=0\csname karo@h@#1\endcsname\relax#2\else#3\fi
  \else#3\fi}
\makeatother
\begin{document}

\noindent{\LARGE\bfseries Potenzen}\par\smallskip
\noindent{\small\color{mbgrau}Selbstlernheft Klasse 9 \textperiodcentered{} mit Lösungsheft \textperiodcentered{} Kennung VUD}\par\medskip
\noindent\textbf{So nutzt du das Heft:} Lies im Abschnitt das Beispiel. Rechne dann die Aufgaben der Reihe nach und vergleiche jede sofort mit dem Lösungsheft. Kreuze „kann ich“ an, wenn du die letzte Aufgabe (\emph{Ziel}) allein schaffst.\par
\noindent Mehr Aufgaben zu einem Abschnitt: Kennung deinem Nachhilfelehrer schicken (z.\,B. „VUD mehr B“).\par\medskip
{\arrayrulecolor{mbgrau}\renewcommand{\arraystretch}{1.3}
\noindent\begin{tabular}{|>{\bfseries}p{10mm}|p{112mm}|>{\raggedleft\arraybackslash}p{10mm}|>{\centering\arraybackslash}p{14mm}|}\hline
\rowcolor{black!8} & \textbf{Abschnitt} & \textbf{Seite} & \textbf{kann ich}\\\hline
A & Potenz als Malkette & \pageref{s:A} & $\square$\\\hline
B & Minus und Komma in der Potenz & \pageref{s:B} & $\square$\\\hline
C & Vergleichen und Exponent finden & \pageref{s:C} & $\square$\\\hline
D & Negative Hochzahl und hoch null & \pageref{s:D} & $\square$\\\hline
T & Probetest & \pageref{s:T} & $\square$\\\hline
\end{tabular}}\par\bigskip

\begin{tcolorbox}[colback=mbkasten,colframe=mbgrau,boxrule=0.5pt,arc=1mm,left=3mm,right=3mm,top=2mm,bottom=2mm,title={\bfseries Formel auf einen Blick},coltitle=black,colbacktitle=black!10]
{\Large $a^n = \underbrace{a \cdot a \cdot \ldots \cdot a}_{n \text{ Faktoren}}$ \qquad $a^{-n} = \dfrac{1}{a^n}$ \qquad $a^0 = 1$}\par\medskip
\textbf{So gehst du vor}\par
\nr{1}\textbf{A} Potenz ausrechnen: Malkette hinschreiben, Schritt für Schritt malnehmen\par
\nr{2}\textbf{B} Minus in Klammern: gerade Hochzahl gibt plus, ungerade minus. Ohne Klammer bleibt das Minus vor dem Ergebnis\par
\nr{3}\textbf{B} Dezimalzahl: Kommastellen der Faktoren zusammenzählen. Bruch: Zähler und Nenner einzeln hoch nehmen\par
\nr{4}\textbf{C} Vergleichen: erst beide Potenzen ausrechnen, dann das Zeichen setzen\par
\nr{5}\textbf{C} Exponent finden: Basis malnehmen, bis der Wert dasteht, Faktoren zählen\par
\nr{6}\textbf{D} Minus im Exponenten: erst Bruch, dann Dezimalzahl\par
\medskip
\end{tcolorbox}
\medskip\noindent\textbf{Typische Fehler – so nicht}\par\smallskip
\noindent\nr{$\times$}$6^3 = 18$ gerechnet (Basis mal Exponent). Richtig: $6^3 = 6 \cdot 6 \cdot 6 = 216$.\par
\noindent\nr{$\times$}$0{,}3^2 = 0{,}6$ gerechnet (verdoppelt). Richtig: $0{,}3 \cdot 0{,}3 = 0{,}09$.\par
\noindent\nr{$\times$}$5^{-2} = -25$ gerechnet. Richtig: $5^{-2} = \frac{1}{25}$, ein Bruch, keine negative Zahl.\par
\noindent\nr{$\times$}Bei $3^x = 243$ geteilt: $243 : 3 = 81$. Richtig: Faktoren zählen, $x = 5$.\par
\clearpage\label{s:A}\noindent\colorbox{black!12}{\makebox[10mm]{\Large\bfseries\strut A}}\hspace{3mm}{\Large\bfseries Potenz als Malkette}\par\smallskip
\noindent Eine Potenz ist eine Malkette aus gleichen Faktoren. Die \textbf{Basis} unten wird malgenommen, der \textbf{Exponent} oben (die Hochzahl) zählt, wie oft. Wird in einer Sache immer wieder mit derselben Zahl malgenommen (verdoppeln heißt mal 2), ist das eine Potenz.\par\medskip
\noindent\textbf{Formel}\par\nopagebreak\noindent\hspace*{4mm}\begin{minipage}{\dimexpr\linewidth-4mm\relax}$a^n = \underbrace{a \cdot a \cdot \ldots \cdot a}_{n \text{ Faktoren}}$ \qquad {\small $a^2$ Quadratzahl, \ $a^3$ Kubikzahl}\end{minipage}\par\medskip
\noindent\textbf{Beispiel}\quad Schreibe $6^3$ als Malkette und rechne aus. Schreibe $5 \cdot 5 \cdot 5 \cdot 5$ als Potenz und rechne aus.\par\smallskip
\noindent\par\noindent\hangindent6mm\hangafter1\makebox[6mm][l]{\textbf{1}}Unten steht die Basis $6$, oben der Exponent $3$. Der Exponent sagt: drei Faktoren $6$. \quad $6^3 = 6 \cdot 6 \cdot 6$\par\smallskip\par\noindent\hangindent6mm\hangafter1\makebox[6mm][l]{\textbf{2}}Schritt für Schritt malnehmen: \quad $6 \cdot 6 = 36$, \quad $36 \cdot 6 = 216$. \ Also $6^3 = 216$.\par\smallskip\par\noindent\hangindent6mm\hangafter1\makebox[6mm][l]{\textbf{3}}Rückwärts: In $5 \cdot 5 \cdot 5 \cdot 5$ steht viermal der Faktor $5$. Basis $5$, Exponent $4$: \quad $5^4$\par\smallskip\par\noindent\hangindent6mm\hangafter1\makebox[6mm][l]{\textbf{4}}Ausrechnen: \quad $5 \cdot 5 = 25$, \ $25 \cdot 5 = 125$, \ $125 \cdot 5 = 625$. \ Also $5^4 = 625$.\par\smallskip\par\medskip
\noindent\textbf{Aufgaben}\hspace{1em}{\small\color{mbgrau}\karohinweis{A}{Rechne unten im Karo.}{Rechne auf einem extra Blatt.} Vergleiche jede Aufgabe gleich mit dem Lösungsheft.}\par\smallskip
\aufg{\hangindent7mm\hangafter1\nr{1.}Schreibe als Malkette und rechne aus: \ a)~$2^5$ \qquad b)~$10^4$ \\ Schreibe als Potenz und rechne aus: \ c)~$7 \cdot 7 \cdot 7$ \qquad d)~$3 \cdot 3 \cdot 3 \cdot 3 \cdot 3$}{}{}
\medskip
\aufg{\hangindent7mm\hangafter1\nr{2.}Berechne beides. \ a)~$12^2$ und $12 \cdot 2$ \qquad b)~$2^6$ und $2 \cdot 6$ \qquad c)~$5^3$ und $5 \cdot 3$ \qquad d)~$15^2$ und $15 \cdot 2$}{}{}
\medskip
\aufg{\hangindent7mm\hangafter1\nr{3.}Ein Bakterium teilt sich jede Stunde in zwei. \ a)~Wie viele Bakterien sind es nach 1, 2 und 3 Stunden? \ b)~Schreibe die Anzahl nach 10 Stunden als Potenz und rechne sie aus.}{}{}
\medskip
\aufg{\hangindent7mm\hangafter1\nr{4.}Ein Blatt Papier ist $0{,}1$\,mm dick. Du faltest es in der Mitte, dann noch einmal, insgesamt 7-mal. Jedes Falten verdoppelt die Zahl der Lagen. Ist der Stapel dann dicker als 10\,mm?}{}{{\scriptsize\color{mbgrau}Ziel\hspace{2mm}}}
\medskip
\par\vspace{3mm}\setlength{\kr}{\dimexpr\pagegoal-\pagetotal-10mm\relax}%
\ifdim\kr>12mm\karomerk{k}{A}\noindent{\scriptsize\color{mbgrau}Platz zum Rechnen}\par\nointerlineskip\vspace{1mm}\noindent
\begin{tikzpicture}\clip (0,0) rectangle ({\linewidth-0.5mm},\kr);\draw[black!16,line width=0.3pt,step=5mm] (0,0) grid (\linewidth,\kr);\end{tikzpicture}\fi\par
\clearpage\label{s:B}\noindent\colorbox{black!12}{\makebox[10mm]{\Large\bfseries\strut B}}\hspace{3mm}{\Large\bfseries Minus und Komma in der Potenz}\par\smallskip
\noindent Steht eine negative Zahl in Klammern, wird sie mit ihrem Minus malgenommen. Ohne Klammer gehört der Exponent nur zur Zahl; das Minus bleibt davor.\par\medskip
\noindent\textbf{Formel}\par\nopagebreak\noindent\hspace*{4mm}\begin{minipage}{\dimexpr\linewidth-4mm\relax}$(-a)^n$: \ $n$ gerade $\to$ plus, \ $n$ ungerade $\to$ minus \qquad $-a^n = -(a^n)$\par $\left(\dfrac{p}{q}\right)^n = \dfrac{p^n}{q^n}$ \qquad {\small Dezimalzahl: Kommastellen zusammenzählen}\end{minipage}\par\medskip
\noindent\textbf{Beispiel}\quad Berechne \ a)~$(-5)^2$ \qquad b)~$(-4)^3$ \qquad c)~$-5^2$ \qquad d)~$0{,}3^2$ \qquad e)~$\left(\frac{3}{4}\right)^2$\par\smallskip
\noindent\par\noindent\hangindent6mm\hangafter1\makebox[6mm][l]{\textbf{1}}Klammer: Die ganze Zahl $-5$ wird malgenommen. Minus mal Minus gibt Plus: \quad $(-5)^2 = (-5) \cdot (-5) = 25$\par\smallskip\par\noindent\hangindent6mm\hangafter1\makebox[6mm][l]{\textbf{2}}Drei Faktoren mit Minus: \quad $(-4) \cdot (-4) \cdot (-4) = 16 \cdot (-4) = -64$\par\smallskip\par\noindent\hangindent6mm\hangafter1\makebox[6mm][l]{\textbf{3}}Keine Klammer: Der Exponent gehört nur zur $5$, das Minus bleibt davor: \quad $-5^2 = -(5 \cdot 5) = -25$\par\smallskip\par\noindent\hangindent6mm\hangafter1\makebox[6mm][l]{\textbf{4}}Dezimalzahl: \quad $0{,}3 \cdot 0{,}3$; \ $3 \cdot 3 = 9$, Kommastellen $1 + 1 = 2$: \quad $0{,}3^2 = 0{,}09$\par\smallskip\par\noindent\hangindent6mm\hangafter1\makebox[6mm][l]{\textbf{5}}Bruch: Zähler mal Zähler, Nenner mal Nenner: \quad $\left(\frac{3}{4}\right)^2 = \frac{3 \cdot 3}{4 \cdot 4} = \frac{9}{16}$\par\smallskip\par\medskip
\noindent\textbf{Aufgaben}\hspace{1em}{\small\color{mbgrau}\karohinweis{B}{Rechne unten im Karo.}{Rechne auf einem extra Blatt.} Vergleiche jede Aufgabe gleich mit dem Lösungsheft.}\par\smallskip
\aufg{\hangindent7mm\hangafter1\nr{1.}Berechne \ a)~$(-6)^2$ \qquad b)~$(-2)^5$ \qquad c)~$(-1)^8$ \qquad d)~$(-10)^3$}{}{}
\medskip
\aufg{\hangindent7mm\hangafter1\nr{2.}Berechne beides. \ a)~$(-7)^2$ und $-7^2$ \qquad b)~$(-1)^4$ und $-1^4$ \qquad c)~$(-3)^3$ und $-3^3$}{}{}
\medskip
\aufg{\hangindent7mm\hangafter1\nr{3.}Berechne \ a)~$0{,}5^2$ \qquad b)~$1{,}2^2$ \qquad c)~$0{,}2^3$ \qquad d)~$\left(\frac{2}{5}\right)^2$ \qquad e)~$\left(\frac{1}{2}\right)^3$}{}{}
\medskip
\aufg{\hangindent7mm\hangafter1\nr{4.}Hier stehen vier Zahlen. Unterstreiche die kleinste. \quad $-6^2$ \qquad $(-6)^2$ \qquad $(-3)^3$ \qquad $-0{,}5^2$}{}{{\scriptsize\color{mbgrau}Ziel\hspace{2mm}}}
\medskip
\par\vspace{3mm}\setlength{\kr}{\dimexpr\pagegoal-\pagetotal-10mm\relax}%
\ifdim\kr>12mm\karomerk{k}{B}\noindent{\scriptsize\color{mbgrau}Platz zum Rechnen}\par\nointerlineskip\vspace{1mm}\noindent
\begin{tikzpicture}\clip (0,0) rectangle ({\linewidth-0.5mm},\kr);\draw[black!16,line width=0.3pt,step=5mm] (0,0) grid (\linewidth,\kr);\end{tikzpicture}\fi\par
\clearpage\label{s:C}\noindent\colorbox{black!12}{\makebox[10mm]{\Large\bfseries\strut C}}\hspace{3mm}{\Large\bfseries Vergleichen und Exponent finden}\par\smallskip
\noindent Welche Potenz größer ist, sieht man nicht an der Basis oder am Exponenten – rechne beide aus. Ist der Exponent gesucht, nimm die Basis so oft mal, bis der Wert dasteht.\par\medskip
\noindent\textbf{Formel}\par\nopagebreak\noindent\hspace*{4mm}\begin{minipage}{\dimexpr\linewidth-4mm\relax}$a^x = b$: \ $a,\ a \cdot a,\ a \cdot a \cdot a,\ \ldots$ bis $b$ \ $\to$ \ Zahl der Faktoren $= x$\end{minipage}\par\medskip
\noindent\textbf{Beispiel}\quad a)~Welche Zahl ist größer, $2^5$ oder $5^2$? \qquad b)~Für welches $x$ gilt $3^x = 243$?\par\smallskip
\noindent\par\noindent\hangindent6mm\hangafter1\makebox[6mm][l]{\textbf{1}}Beide ausrechnen: \quad $2^5 = 2 \cdot 2 \cdot 2 \cdot 2 \cdot 2 = 32$, \quad $5^2 = 5 \cdot 5 = 25$. \ Also $2^5 > 5^2$.\par\smallskip\par\noindent\hangindent6mm\hangafter1\makebox[6mm][l]{\textbf{2}}Exponent finden: Die $3$ immer wieder malnehmen, bis $243$ dasteht: \quad $3,\ 9,\ 27,\ 81,\ 243$\par\smallskip\par\noindent\hangindent6mm\hangafter1\makebox[6mm][l]{\textbf{3}}Faktoren zählen: $243 = 3 \cdot 3 \cdot 3 \cdot 3 \cdot 3$, das sind fünf. \ Also $x = 5$.\par\smallskip\par\noindent\hangindent6mm\hangafter1\makebox[6mm][l]{\textbf{4}}Probe: \quad $3^5 = 243$ \checkmark\par\smallskip\par\medskip
\noindent\textbf{Aufgaben}\hspace{1em}{\small\color{mbgrau}\karohinweis{C}{Rechne unten im Karo.}{Rechne auf einem extra Blatt.} Vergleiche jede Aufgabe gleich mit dem Lösungsheft.}\par\smallskip
\aufg{\hangindent7mm\hangafter1\nr{1.}Setze $<$, $=$ oder $>$ ein. \ a)~$3^2 \;\square\; 2^3$ \qquad b)~$10^2 \;\square\; 2^7$ \qquad c)~$4^2 \;\square\; 2^4$ \qquad d)~$1^{10} \;\square\; 10^1$}{}{}
\medskip
\aufg{\hangindent7mm\hangafter1\nr{2.}Für welches $x$ gilt die Gleichung? \ a)~$2^x = 64$ \qquad b)~$5^x = 625$ \qquad c)~$10^x = 100\,000$ \qquad d)~$3^x = 27$}{}{}
\medskip
\aufg{\hangindent7mm\hangafter1\nr{3.}Hier stehen drei Zahlen. Unterstreiche die größte. \quad $5^3$ \qquad $2^7$ \qquad $11^2$}{}{}
\medskip
\aufg{\hangindent7mm\hangafter1\nr{4.}Auf einem Teich wachsen Algen. Heute bedecken sie 1\,m$^2$, jeden Tag verdoppelt sich die Fläche. Der Teich ist 512\,m$^2$ groß. Ist er in einer Woche ganz zugewachsen?}{}{{\scriptsize\color{mbgrau}Ziel\hspace{2mm}}}
\medskip
\par\vspace{3mm}\setlength{\kr}{\dimexpr\pagegoal-\pagetotal-10mm\relax}%
\ifdim\kr>12mm\karomerk{k}{C}\noindent{\scriptsize\color{mbgrau}Platz zum Rechnen}\par\nointerlineskip\vspace{1mm}\noindent
\begin{tikzpicture}\clip (0,0) rectangle ({\linewidth-0.5mm},\kr);\draw[black!16,line width=0.3pt,step=5mm] (0,0) grid (\linewidth,\kr);\end{tikzpicture}\fi\par
\clearpage\label{s:D}\noindent\colorbox{black!12}{\makebox[10mm]{\Large\bfseries\strut D}}\hspace{3mm}{\Large\bfseries Negative Hochzahl und hoch null}\par\smallskip
\noindent Ein Minus im Exponenten macht keine negative Zahl, sondern einen Bruch. Jede Zahl außer $0$ hoch $0$ ist $1$.\par\medskip
\noindent\textbf{Formel}\par\nopagebreak\noindent\hspace*{4mm}\begin{minipage}{\dimexpr\linewidth-4mm\relax}$a^{-n} = \dfrac{1}{a^n}$ \qquad $a^0 = 1$\end{minipage}\par\medskip
\noindent\textbf{Beispiel}\quad Schreibe $5^{-2}$ als Bruch und als Dezimalzahl. Ist $5^{-2}$ größer oder kleiner als $0{,}05$? Berechne auch $9^0$.\par\smallskip
\noindent\par\noindent\hangindent6mm\hangafter1\makebox[6mm][l]{\textbf{1}}Minus in der Hochzahl heißt „1 geteilt durch“: \quad $5^{-2} = \frac{1}{5^2}$\par\smallskip\par\noindent\hangindent6mm\hangafter1\makebox[6mm][l]{\textbf{2}}Unten ausrechnen: $5^2 = 25$, \ also \ $5^{-2} = \frac{1}{25}$\par\smallskip\par\noindent\hangindent6mm\hangafter1\makebox[6mm][l]{\textbf{3}}Als Dezimalzahl: erweitern auf Hundertstel: \quad $\frac{1}{25} = \frac{4}{100} = 0{,}04$ \quad (oder $1 : 25$). Geht Erweitern nicht: schriftlich teilen, etwa $1 : 8 = 0{,}125$.\par\smallskip\par\noindent\hangindent6mm\hangafter1\makebox[6mm][l]{\textbf{4}}Vergleichen: \quad $0{,}04 < 0{,}05$, \ also \ $5^{-2} < 0{,}05$\par\smallskip\par\noindent\hangindent6mm\hangafter1\makebox[6mm][l]{\textbf{5}}Hoch null: $9^2 = 81$, \ $9^1 = 9$, \ $9^0 = 9 : 9 = 1$ \ (jede Stufe tiefer: durch $9$ teilen).\par\smallskip\par\medskip
\noindent\textbf{Aufgaben}\hspace{1em}{\small\color{mbgrau}\karohinweis{D}{Rechne unten im Karo.}{Rechne auf einem extra Blatt.} Vergleiche jede Aufgabe gleich mit dem Lösungsheft.}\par\smallskip
\aufg{\hangindent7mm\hangafter1\nr{1.}Schreibe als Bruch. \ a)~$3^{-2}$ \qquad b)~$10^{-3}$ \qquad c)~$2^{-4}$ \qquad d)~$7^{-1}$ \qquad e)~$6^0$}{}{}
\medskip
\aufg{\hangindent7mm\hangafter1\nr{2.}Schreibe als Bruch und als Dezimalzahl. \ a)~$2^{-1}$ \qquad b)~$5^{-1}$ \qquad c)~$10^{-2}$ \qquad d)~$4^{-1}$}{}{}
\medskip
\aufg{\hangindent7mm\hangafter1\nr{3.}Setze $<$, $=$ oder $>$ ein. \ a)~$8^{-1} \;\square\; 0{,}1$ \qquad b)~$10^{-3} \;\square\; 0{,}01$ \qquad c)~$2^{-2} \;\square\; 0{,}2$ \qquad d)~$3^0 \;\square\; 0{,}9$}{}{}
\medskip
\aufg{\hangindent7mm\hangafter1\nr{4.}Ordne die Zahlen der Größe nach. Beginne mit der kleinsten: \quad $5^{-1}$ \qquad $0{,}15$ \qquad $20^{-1}$ \qquad $0{,}3$ \qquad $4^{-1}$}{}{{\scriptsize\color{mbgrau}Ziel\hspace{2mm}}}
\medskip
\par\vspace{3mm}\setlength{\kr}{\dimexpr\pagegoal-\pagetotal-10mm\relax}%
\ifdim\kr>12mm\karomerk{k}{D}\noindent{\scriptsize\color{mbgrau}Platz zum Rechnen}\par\nointerlineskip\vspace{1mm}\noindent
\begin{tikzpicture}\clip (0,0) rectangle ({\linewidth-0.5mm},\kr);\draw[black!16,line width=0.3pt,step=5mm] (0,0) grid (\linewidth,\kr);\end{tikzpicture}\fi\par
\clearpage\label{s:T}\noindent\colorbox{black!12}{\makebox[10mm]{\Large\bfseries\strut T}}\hspace{3mm}{\Large\bfseries Probetest}\par\smallskip
\noindent Gemischt wie in der Prüfung, ohne Beispiel. Ohne Taschenrechner. Etwa 15 Minuten.\par\medskip
\noindent\textbf{Aufgaben}\hspace{1em}{\small\color{mbgrau}\karohinweis{T}{Rechne unten im Karo.}{Rechne auf einem extra Blatt.} Vergleiche jede Aufgabe gleich mit dem Lösungsheft.}\par\smallskip
\aufg{\hangindent7mm\hangafter1\nr{1.}Für welches $x$ gilt $8^x = 4096$?}{}{}
\medskip
\aufg{\hangindent7mm\hangafter1\nr{2.}Berechne \ a)~$(-6)^3$ \qquad b)~$-6^2$ \qquad c)~$0{,}4^3$ \qquad d)~$(-5)^0$ \qquad e)~$10^{-1}$}{}{}
\medskip
\aufg{\hangindent7mm\hangafter1\nr{3.}Hier stehen drei Zahlen. Unterstreiche die größte. \quad $4^4$ \qquad $3^5$ \qquad $15^2$}{}{}
\medskip
\aufg{\hangindent7mm\hangafter1\nr{4.}Setze $<$, $=$ oder $>$ ein. \qquad $2^{-4} \;\square\; 0{,}06$}{}{}
\medskip
\aufg{\hangindent7mm\hangafter1\nr{5.}Ein Video wird geteilt: Wer es bekommt, schickt es nach einer Stunde an 4 neue Personen. Am Anfang hat es eine Person. Nach 1 Stunde bekommen es 4 Personen, nach 2 Stunden 16 neue. Wie viele bekommen es nach 5 Stunden neu – mehr als 1000?}{}{{\scriptsize\color{mbgrau}Ziel\hspace{2mm}}}
\medskip
\par\vspace{3mm}\setlength{\kr}{\dimexpr\pagegoal-\pagetotal-10mm\relax}%
\ifdim\kr>12mm\karomerk{k}{T}\noindent{\scriptsize\color{mbgrau}Platz zum Rechnen}\par\nointerlineskip\vspace{1mm}\noindent
\begin{tikzpicture}\clip (0,0) rectangle ({\linewidth-0.5mm},\kr);\draw[black!16,line width=0.3pt,step=5mm] (0,0) grid (\linewidth,\kr);\end{tikzpicture}\fi\par
\end{document}
